\textbf{结论名称：}构造等比数列（\(a_{n+1}=pa_n+q n+r\) 型）

\textbf{适用条件：}递推关系形如 \(a_{n+1}=p a_n+q n+r\)，其中 \(p\neq0,1\)，\(q,r\) 为常数（\(q\neq0\) 时存在一次项）。

\textbf{变量说明：}\(p,q,r\) 为已知常数；\(A=\frac{q}{p-1}\)，\(B=\frac{r}{p-1}+\frac{q}{(p-1)^2}\) 为待定系数；设辅助数列 \(b_n=a_n+A n+B\)。

\textbf{核心公式：}
\[
\boxed{b_{n+1}=p\,b_n},\quad \text{其中 } b_n=a_n+\frac{q}{p-1}n+\frac{r}{p-1}+\frac{q}{(p-1)^2}.
\]
通项公式为
\[
a_n=b_1 p^{\,n-1}-\frac{q}{p-1}n-\frac{r}{p-1}-\frac{q}{(p-1)^2}.
\]

\textbf{使用场景：}处理含有一次项 \(q n+r\) 的线性递推求通项，用待定系数法消去 \(n\) 的显式依赖，将其转化为等比数列模型。

\textbf{简单例子：}已知 \(a_{n+1}=2a_n+3n+1\)，且 \(a_1=1\)。  
由 \(p=2,q=3,r=1\) 得 \(A=\frac{3}{2-1}=3\)，\(B=\frac{1}{1}+\frac{3}{1^2}=4\)。  
令 \(b_n=a_n+3n+4\)，则 \(b_{n+1}=2b_n\)，\(b_1=1+3+4=8\)。  
故 \(b_n=8\cdot2^{\,n-1}=2^{\,n+2}\)，通项 \(a_n=2^{\,n+2}-3n-4\)。  
验证：\(a_2=2\cdot1+3\cdot1+1=6\)，由公式得 \(2^4-6-4=6\)，一致。